\providecommand{\pdfxopt}{a-1b,mathxmp}

%\documentclass[10pt]{article}
%\usepackage{beamerarticle}
\documentclass[ignorenonframetext,slidestop]{beamer}

\input{/home/asreeve/current/class/ers102/pre_latex/luatex_pre.tex}

\setbeamertemplate{enumerate item}[square]
\setbeamertemplate{enumerate subitem}[circle]
\setbeamertemplate{enumerate subsubitem}[square]
\setbeamercolor*{enumerate item}{fg=red,bg=blue}
\setbeamercolor*{enumerate subitem}{fg=blue,bg=red}

\begin{document}

\begin{frame}{Metamorphic Rocks}
  \begin{enumerate}
  \item Metamorphic Rocks are rocks that have formed from other rocks.
  \item The precursor (original) rock refereed to as parent material. Parent material impacts mineralogy of rock. Parent material of slate is shale (very fine grained (w/ clay minerals) sedimentary rock).
  \item Heat and pressure transforms or metamorphoses other rocks into this rock group. 
  \item Partial melting is possible. Minerals are rearranged and altered in metamorphic rocks.
  \item What you learned about minerals in the igneous rock lab is applicable to this lab.
  \end{enumerate}
\end{frame}

\begin{frame}{Metamorphic Rock Classification}
  Primary metamorphic rock classification is based on texture with further classification based on mineralogy.
  \begin{enumerate}
  \item Foliated (layered fabric)
    \begin{itemize}
    \item Shale
    \item Phyllite
    \item Schist
    \item Gneiss
    \end{itemize}
  \item Non-foliated (massive texture)
    \begin{itemize}
    \item Marble (reacts with acid! Soft.)
    \item Quartzite (really hard)
    \end{itemize}
  \end{enumerate}
\end{frame}

\begin{frame}{Metamorphic Minerals and Grade}
    Metamorphism results in some minerals that usually occur (in many cases only occur) in metamorphic rocks. The minerals that form and the rock's texture, indicate the conditions the rock formed in. What was the temperature and pressure of formation? How deep in the earth was the rock created?
    
    \begin{enumerate}
    \item chlorite: a low grade mineral, forms under relatively low T and P
    \item staurolite: an intermediate grade mineral, forms under moderate T and P conditions
    \item garnet: intermediate to high grade mineral (without complete melting of rock!)
    \end{enumerate}

    Similarly:
    \begin{enumerate}
    \item slate: typically low grade metamorphic rock, minerals have been reoriented
    \item gneiss: typically high grade rock, partial melting allowed minerals to separate into bands 
    \end{enumerate}
\end{frame}


\begin{frame}{Doing the Lab}{Mineral ID}
  \begin{enumerate}
  \item Minerals:numbers, rocks:letters
  \item Watch the videos or examine hand samples and fill in the attached tables.
  \item More rows on the table than there are samples.
  \item Table of minerals contains minerals not in this lab, they appear in other labs
  \item Once properties defined, ID mineral sample and fill in table
  \item one new test appears in the videos, placing a drop of acid on a sample. If a sample contains calcite or similar mineral, it will fizz or effervesce.
  \end{enumerate}
The following minerals are in this lab: chlorite, garnet, staurolite, kyanite, muscovite, biotite. 
\end{frame}

\begin{frame}{Doing the Lab}{Rock ID}
  \begin{enumerate}
  \item group rocks by texture
  \item compare with rock descriptions and ID rocks
  \end{enumerate}
The following rock samples are present: slate, phyllite, schist, gneiss, hornfels, marble, quartzite. 
\end{frame}
\end{document}
